%%% ---------------------------------------------------------------------------
%%% setup.tex —— 报告元数据
%%%
%%% 只填内容，不写格式。三种 type 各用其中一组，用不到的留空即可
%%% （留空的字段不会出现在封面上）。
%%% ---------------------------------------------------------------------------

%%% ---- 通用 ----
\title{不同优化器在 CIFAR-10 上的收敛性对比}
% \subtitle{副标题，不需要就留空}
\school{中国科学技术大学}
\department{信息科学技术学院}
% \logo{\includegraphics[height=2.2cm]{logo}}   % 校徽，图放在 figures/

%%% ---- type = lab / course 用这一组 ----
\course{深度学习实验}
\experimentno{三}                    % 不写就不显示「实验三」
\experimentname{优化器对比实验}
\studentname{张三}
\studentid{PB20000001}
\classname{计算机 2001 班}
\partners{李四、王五}                % 独立完成就留空
\instructor{赵老师}
\expdate{2026 年 7 月 20 日}
\submitdate{2026 年 7 月 26 日}

%%% ---- type = tech 用这一组（GB/T 7713.3—2014）----
% \reportno{USTC-TR-2026-017}
% \seclevel{内部}
% \organization{中国科学技术大学 信息科学技术学院}
% \authors{张三、李四}
% \reviewer{赵老师}
% \approver{钱主任}
% \reportversion{V1.0}
% \keywords{优化器；收敛性；CIFAR-10}   % 技术报告的摘要块会用到

%%% ---- 页眉短标题 ----
\runningtitle{优化器对比实验}

%%% ---- 参考文献数据库 ----
%%% refs.bib 是指向 ../common/refs.bib 的符号链接，三套模板共用同一份文献库。
\addbibresource{refs.bib}
